(1) 令未知输入为比特1、共享对为2和3。定义 \(|\Phi^\pm\rangle=(|00\rangle\pm|11\rangle)/\sqrt2\)、\(|\Psi^\pm\rangle=(|01\rangle\pm|10\rangle)/\sqrt2\)。逐项展开：
\[
|\psi\rangle_1|\Phi^+\rangle_{23}=\tfrac12[
|\Phi^+\rangle_{12}|\psi\rangle_3+
|\Phi^-\rangle_{12}Z|\psi\rangle_3+
|\Psi^+\rangle_{12}X|\psi\rangle_3+
|\Psi^-\rangle_{12}XZ|\psi\rangle_3].
\]
四个Bob条件态范数均为1，因此四种Bell结果概率各 \(1/4\)，与未知 \(\alpha,\beta\) 无关。

(2) Alice通过经典信道发送结果标签，Bob依次施加 \(I,Z,X,ZX\)，每种情形均把对应Pauli误差逆转（至多差一个全局相位）。未收到经典消息时，Bob的四分支平均约化态为 \(I_2/2\)，不含未知输入态的可读取信息，因此纠缠本身不允许超光速传信。

(3) 传态映射对资源态线性。理想Bell资源给恒等信道，完全混合资源 \(I_4/4\) 给完全退极化输出 \(I_2/2\)，因此
\[
\boxed{\mathcal T_p(\rho)=p\rho+(1-p)I_2/2}.
\]
任意纯输入 \(\rho=|\psi\rangle\langle\psi|\) 的保真度为
\(F=\langle\psi|\mathcal T_p(\rho)|\psi\rangle=p+(1-p)/2=(1+p)/2\)，与Bloch球方向无关，故均匀平均仍为同值。

(4) 对未知纯单比特，最优纯经典测量—重制备协议的平均保真度上界为 \(2/3\)，所以
\[
\boxed{F_{\rm avg}>2/3\iff p>1/3}.
\]
\(p=1/3\) 仅与经典界持平，不能宣称严格优势；此处亦为Werner资源的纠缠阈值。Bob需要Alice的两比特经典信息，且只后选择某一Bell结果会把成功率降低至 \(1/4\)，不能从资源账本中删去失败试验。
